% ==============================================================================
% Chapter 4: Risk, Return, and Portfolio Mathematics
% Course: MSQF 532 Financial Mathematics
% Continuous Notes Architecture: Lectures from 28/10/2026 to 30/10/2026
% ==============================================================================

\chapter{Risk, Return, and Portfolio Mathematics}
\label{ch:risk_return_portfolio_mathematics}

This chapter establishes the statistical, probabilistic, and quantitative foundations of investment analysis, asset pricing, and portfolio theory in accordance with Unit~IV of the syllabus. We examine the fundamental trade-off between financial risk and return, explore classical measures of central tendency (arithmetic, geometric, and harmonic means), formalize absolute and relative dispersion techniques (variance, standard deviation, and coefficient of variation), analyze bivariate market co-movement through Pearson correlation, and introduce the quantitative foundations of foreign exchange exposure.

% ==============================================================================
% 4.1 Financial Risk and Return Axioms & Measures of Central Tendency
% ==============================================================================
\section{Financial Risk and Return Axioms \& Measures of Central Tendency}
\lecturedate{28/10/2026}

\begin{definition}{The Risk--Return Spectrum in Quantitative Finance}{def_risk_return_spectrum}
In financial mathematics and portfolio theory, investment decision-making is governed by two complementary economic dimensions:
\begin{enumerate}[leftmargin=*, topsep=2pt, itemsep=3pt]
    \item \textbf{Financial Risk (Uncertainty and Dispersion):} The possibility that the realized return on an investment deviates adversely from its anticipated or expected value. It encompasses capital impairment, cash flow volatility, default risk, and market downside.
    \item \textbf{Financial Return (Potential Profit):} The net economic reward, yield, or capital gain generated by committing economic capital over a designated temporal horizon.
\end{enumerate}
The foundational axiom of modern financial economics asserts that risk and expected return are positively correlated: rational, risk-averse economic agents demand higher expected compensation for absorbing elevated levels of variance or downside uncertainty.
\end{definition}

\subsection{Measures of Central Tendency for Financial Return Series}
An \textbf{average} is a single representative central value summarizing the location of an entire distribution of financial data points. In portfolio analysis and performance attribution, five primary measures of central tendency are employed:

\begin{formulabox}[Mathematical Measures of Central Tendency]
\begin{itemize}[leftmargin=*, noitemsep]
    \item \textbf{Simple Arithmetic Mean ($\bar{X}$):}
    The arithmetic mean sums all observed outcomes and divides by the sample cardinality $N$:
    \begin{equation}
        \bar{X} = \frac{1}{N}\sum_{i=1}^{N} X_i = \frac{X_1 + X_2 + \cdots + X_N}{N}
    \end{equation}

    \item \textbf{Weighted Arithmetic Mean ($\bar{X}_w$):}
    Accounts for differential economic importance, asset market capitalization, or portfolio allocation weights $w_i$:
    \begin{equation}
        \bar{X}_w = \frac{\sum_{i=1}^{N} w_i X_i}{\sum_{i=1}^{N} w_i} \qquad \text{with} \quad \sum_{i=1}^{N} w_i = 1 \implies \bar{X}_w = \sum_{i=1}^{N} w_i X_i
    \end{equation}

    \item \textbf{Mode (Modal Value):}
    The observation that occurs with the highest empirical frequency in a data series:
    \begin{equation}
        \text{Mode} = \arg\max_{x} f(x)
    \end{equation}

    \item \textbf{Geometric Mean ($\GM$):}
    The $n$-th root of the product of $n$ positive financial observations or growth multipliers:
    \begin{equation}
        \GM = \sqrt[n]{\prod_{i=1}^{n} X_i} = \left(X_1 \cdot X_2 \cdots X_n\right)^{1/n} = \exp\left(\frac{1}{n}\sum_{i=1}^{n} \ln X_i\right)
    \end{equation}
    The geometric mean is the strictly appropriate metric for compounding multi-period investment returns and evaluating the Compound Annual Growth Rate (CAGR).

    \item \textbf{Harmonic Mean ($\HM$):}
    The reciprocal of the arithmetic mean of the reciprocals of the data elements:
    \begin{equation}
        \HM = \frac{n}{\sum_{i=1}^{n} \frac{1}{X_i}} = \frac{n}{\frac{1}{X_1} + \frac{1}{X_2} + \cdots + \frac{1}{X_n}}
    \end{equation}
    The harmonic mean is utilized in finance when averaging rates, price-to-earnings ($P/E$) multiples, and executing fixed-rupee cost averaging (SIP investment strategies).
\end{itemize}
\begin{equation*}
    \textbf{Fundamental Classical Inequality:} \quad \AM \ge \GM \ge \HM \quad (\text{for all positive } X_i > 0)
\end{equation*}
Equality holds if and only if all observed elements are identical ($X_1 = X_2 = \cdots = X_n$).
\end{formulabox}

\lecturedate{28/10/2026}
\begin{example}{Arithmetic Mean of Monthly Employee Incomes}{ex_am_incomes}
\textbf{Problem:} \\
The payroll audit records of a commercial financial enterprise disclose the following monthly incomes (in \inr) for 10 corporate employees:
\begin{center}
    $14{,}780, \quad 15{,}760, \quad 26{,}690, \quad 27{,}750, \quad 24{,}840,$ \\
    $24{,}920, \quad 16{,}100, \quad 17{,}810, \quad 27{,}050, \quad 26{,}950$
\end{center}
Calculate the simple arithmetic mean of the monthly income distribution.

\tcbline
\textbf{Solution:} \\
Let $X_i$ denote the monthly income of employee $i$ ($i = 1, 2, \dots, 10$), with sample size $N = 10$.

\begin{center}
\small
\renewcommand{\arraystretch}{1.15}
\begin{tabularx}{0.7\linewidth}{>{\centering\arraybackslash}p{2.5cm} >{\raggedleft\arraybackslash}X}
\toprule
\textbf{Employee ($i$)} & \textbf{Monthly Income ($X_i$ in \inr)} \\
\midrule
1  & 14,780 \\
2  & 15,760 \\
3  & 26,690 \\
4  & 27,750 \\
5  & 24,840 \\
6  & 24,920 \\
7  & 16,100 \\
8  & 17,810 \\
9  & 27,050 \\
10 & 26,950 \\
\midrule
\textbf{Aggregate Sum ($\sum X_i$)} & \textbf{\inr 2,22,650} \\
\bottomrule
\end{tabularx}
\end{center}

Applying the simple arithmetic mean formulation:
\begin{equation*}
    \bar{X} = \frac{\sum_{i=1}^{N} X_i}{N} = \frac{2{,}22{,}650}{10} = \mathbf{\inr 22{,}265}
\end{equation*}
\textit{Economic Interpretation:} The representative central earning baseline across the 10 corporate employees is $\mathbf{\inr 22{,}265}$ per month.
\end{example}

\lecturedate{28/10/2026}
\begin{example}{Modal and Harmonic Mean Valuation of Asset Series}{ex_mode_hm}
\textbf{Problem:} \\
Solve the following two foundational central tendency evaluations:
\begin{enumerate}[label=(\alph*), leftmargin=*]
    \item Determine the mode for the discrete transaction price series:
    \begin{center}
        $3, \quad 5, \quad 8, \quad 5, \quad 4, \quad 5, \quad 9, \quad 3$
    \end{center}
    \item Compute the harmonic mean ($\HM$) for the following series of 8 operational asset rates:
    \begin{center}
        $5, \quad 7, \quad 10, \quad 13, \quad 15, \quad 20, \quad 25, \quad 28$
    \end{center}
\end{enumerate}

\tcbline
\textbf{Solution:} \\
\textbf{Part (a): Modal Value Determination:}
Tabulating the empirical occurrence frequencies of the observation elements:
\begin{center}
\small
\renewcommand{\arraystretch}{1.1}
\begin{tabular}{lcccccccc}
\toprule
\textbf{Transaction Value ($X$)} & 3 & 4 & 5 & 8 & 9 \\
\midrule
\textbf{Frequency ($f$)}         & 2 & 1 & \textbf{3} & 1 & 1 \\
\bottomrule
\end{tabular}
\end{center}
The numerical value $5$ appears with the maximum frequency ($f = 3$). Therefore:
\begin{equation*}
    \mathbf{\text{Mode} = 5}
\end{equation*}

\textbf{Part (b): Harmonic Mean Computation:}
Let $n = 8$. We construct the schedule of reciprocal components $\frac{1}{X_i}$:
\begin{center}
\small
\renewcommand{\arraystretch}{1.12}
\begin{tabularx}{0.75\linewidth}{>{\centering\arraybackslash}p{2.5cm} >{\centering\arraybackslash}p{3.0cm} >{\raggedleft\arraybackslash}X}
\toprule
\textbf{Observation ($i$)} & \textbf{Asset Rate ($X_i$)} & \textbf{Reciprocal Value ($1/X_i$)} \\
\midrule
1 & 5  & $1/5 = 0.20000$ \\
2 & 7  & $1/7 \approx 0.14286$ \\
3 & 10 & $1/10 = 0.10000$ \\
4 & 13 & $1/13 \approx 0.07692$ \\
5 & 15 & $1/15 \approx 0.06667$ \\
6 & 20 & $1/20 = 0.05000$ \\
7 & 25 & $1/25 = 0.04000$ \\
8 & 28 & $1/28 \approx 0.03571$ \\
\midrule
\textbf{Total ($n = 8$)} & --- & $\sum_{i=1}^{8} \frac{1}{X_i} \approx \mathbf{0.71216}$ \\
\bottomrule
\end{tabularx}
\end{center}

Applying the harmonic mean formula:
\begin{equation*}
    \HM = \frac{n}{\sum_{i=1}^{n} \frac{1}{X_i}} = \frac{8}{0.71216} \approx \mathbf{11.233}
\end{equation*}
\textit{Mathematical Verification:} For the identical dataset, $\AM = \frac{123}{8} = 15.375$. Since $\AM (15.375) > \HM (11.233)$, the classical mean inequality $\AM \ge \HM$ is rigorously confirmed.
\end{example}

\subsection{Theoretical Foundations of Portfolio Return and Diversification}
In asset allocation and modern portfolio theory (Markowitz framework), the performance of a multi-asset investment portfolio is structured as follows:
\begin{itemize}[leftmargin=*, topsep=2pt, itemsep=3pt]
    \item \textbf{Portfolio Expected Return ($E[R_p]$):}
    The expected return of a portfolio is the strictly linear weighted average of the expected returns of its constituent security holdings:
    \begin{equation}
        E[R_p] = \sum_{i=1}^{N} w_i E[R_i] \qquad \text{where} \quad \sum_{i=1}^{N} w_i = 1
    \end{equation}
    where $w_i$ represents the fraction of total investment capital committed to security $i$.

    \item \textbf{Diversification Mechanics and Non-Linear Risk Reduction:}
    While portfolio expected return is strictly additive, \textbf{portfolio risk (variance) is sub-additive} whenever the constituent assets are not perfectly positively correlated ($\rho_{ij} < +1$):
    \begin{equation}
        \sigma_p^2 = \sum_{i=1}^{N} w_i^2 \sigma_i^2 + 2\sum_{i=1}^{N}\sum_{j > i} w_i w_j \Cov(R_i, R_j) < \sum_{i=1}^{N} w_i \sigma_i \quad (\text{if } \rho_{ij} < 1)
    \end{equation}
    Diversification pools non-systematic (unsystematic / idiosyncratic) firm-specific shocks, eliminating uncompensated risk without sacrificing expected portfolio yield.
\end{itemize}

% ==============================================================================
% 4.2 Measures of Dispersion, Absolute Risk, and Comparative Stability
% ==============================================================================
\newpage
\section{Measures of Dispersion, Absolute Risk, and Comparative Stability}
\lecturedate{29/10/2026}

\begin{definition}{Absolute versus Relative Dispersion}{def_dispersion_framework}
In financial risk analysis, \textbf{dispersion} quantifies the extent to which asset returns, market prices, or cash flow outcomes scatter or diverge from their central expected mean.
\begin{enumerate}[leftmargin=*, topsep=2pt, itemsep=3pt]
    \item \textbf{Absolute Dispersion:} Expressed in the identical physical or monetary units as the original dataset (e.g. rupees, percentage return points). Primary metrics include Variance ($\sigma^2$) and Standard Deviation ($\sigma$).
    \item \textbf{Relative Dispersion:} Dimensionless, normalized percentage ratios that facilitate direct cross-sectional risk and stability comparisons across assets trading at fundamentally disparate price scales or denominated in different currencies. The primary metric is the Coefficient of Variation ($\CV$).
\end{enumerate}
\end{definition}

\subsection{Formulations of Absolute Dispersion: Variance and Standard Deviation}
\begin{formulabox}[Direct and Shortcut Formulations for Standard Deviation]
\begin{itemize}[leftmargin=*, noitemsep]
    \item \textbf{Population Variance ($\sigma^2$) and Standard Deviation ($\sigma$):}
    \begin{equation}
        \sigma^2 = \frac{\sum_{i=1}^{N} (X_i - \bar{X})^2}{N}, \qquad \sigma = \sqrt{\frac{\sum_{i=1}^{N} (X_i - \bar{X})^2}{N}}
    \end{equation}

    \item \textbf{Sample Variance ($s^2$) and Sample Standard Deviation ($s$):}
    Incorporate Bessel's degrees-of-freedom correction ($n - 1$) to eliminate small-sample downward estimation bias:
    \begin{equation}
        s^2 = \frac{\sum_{i=1}^{n} (X_i - \bar{X})^2}{n - 1}, \qquad s = \sqrt{\frac{\sum_{i=1}^{n} (X_i - \bar{X})^2}{n - 1}}
    \end{equation}

    \item \textbf{Assumed Mean Method (Shortcut Technique):}
    To eliminate cumbersome decimal arithmetic when the true mean $\bar{X}$ is non-integer, an arbitrary round constant $A$ is selected as the \textbf{assumed mean}. Defining deviations $d_i = X_i - A$:
    \begin{equation}
        \bar{X} = A + \frac{\sum d_i}{n}
    \end{equation}
    \begin{equation}
        \sigma = \sqrt{\frac{\sum d_i^2}{n} - \left(\frac{\sum d_i}{n}\right)^2}, \qquad s = \sqrt{\frac{\sum d_i^2 - \frac{(\sum d_i)^2}{n}}{n - 1}}
    \end{equation}

    \item \textbf{Shift of Origin and Scaling Method:}
    When financial observations $X_i$ are tightly clustered around a large base value with small fractional increments, we define a transformed variable $y_i = X_i - A$ or $u_i = \frac{X_i - A}{h}$:
    \begin{equation}
        \sigma_X = h \cdot \sigma_u = h \sqrt{\frac{\sum u_i^2}{n} - \left(\frac{\sum u_i}{n}\right)^2}
    \end{equation}
    Variance is mathematically invariant to shifts in origin ($X_i \to X_i - A$), but scales quadratically with changes of scale ($X_i \to h X_i \implies \sigma^2 \to h^2 \sigma^2$).
\end{itemize}
\end{formulabox}

\subsection{Relative Dispersion and Stability: The Coefficient of Variation (\texorpdfstring{$\CV$}{CV})}
To assess which financial asset exhibits greater price stability or lower risk relative to its earnings level, analysts utilize the \textbf{Coefficient of Variation ($\CV$)}:
\begin{equation}
    \CV = \frac{\text{Standard Deviation}}{\text{Mean}} \times 100 = \frac{\sigma}{\bar{X}} \times 100
\end{equation}
\begin{equation*}
    \textbf{Decision Rule: } \quad \text{Lower } \CV \implies \text{Superior Price Stability / Lower Relative Risk}
\end{equation*}

\lecturedate{29/10/2026}
\begin{example}{Standard Deviation Computation via Assumed Mean and Origin Shift}{ex_sd_methods}
\textbf{Problem:} \\
Solve the following two rigorous standard deviation evaluations:
\begin{enumerate}[label=(\alph*), leftmargin=*]
    \item The blood serum cholesterol readings of 10 persons were recorded as follows:
    \begin{center}
        $240, \quad 260, \quad 290, \quad 245, \quad 255, \quad 288, \quad 272, \quad 263, \quad 277, \quad 251 \quad (\text{in mg/dL})$
    \end{center}
    Compute the standard deviation using the \textbf{Assumed Mean Method} with $A = 260$.
    \item Compute the standard deviation for the following 10 tightly clustered decimal security prices:
    \begin{center}
        $240.12, \quad 240.13, \quad 240.15, \quad 240.12, \quad 240.17,$ \\
        $240.16, \quad 240.22, \quad 240.21, \quad 240.18, \quad 240.17$
    \end{center}
\end{enumerate}

\tcbline
\textbf{Solution:} \\
\textbf{Part (a): Assumed Mean Method ($A = 260$):} \\
Let $d_i = X_i - 260$ with $n = 10$. We construct the deviation table:

\begin{center}
\small
\renewcommand{\arraystretch}{1.12}
\begin{tabularx}{0.75\linewidth}{>{\centering\arraybackslash}p{2.0cm} >{\centering\arraybackslash}p{3.0cm} >{\centering\arraybackslash}p{3.0cm} >{\raggedleft\arraybackslash}X}
\toprule
\textbf{Index ($i$)} & \textbf{Reading ($X_i$)} & \textbf{Deviation ($d_i = X_i - 260$)} & \textbf{Squared ($d_i^2$)} \\
\midrule
1  & 240 & $-20$ & 400 \\
2  & 260 & $0$   & 0 \\
3  & 290 & $+30$ & 900 \\
4  & 245 & $-15$ & 225 \\
5  & 255 & $-5$  & 25 \\
6  & 288 & $+28$ & 784 \\
7  & 272 & $+12$ & 144 \\
8  & 263 & $+3$  & 9 \\
9  & 277 & $+17$ & 289 \\
10 & 251 & $-9$  & 81 \\
\midrule
\textbf{Total} & --- & $\sum d_i = \mathbf{41}$ & $\sum d_i^2 = \mathbf{2{,}857}$ \\
\bottomrule
\end{tabularx}
\end{center}

1. Arithmetic Mean:
\begin{equation*}
    \bar{X} = A + \frac{\sum d_i}{n} = 260 + \frac{41}{10} = \mathbf{264.1 \text{ mg/dL}}
\end{equation*}

2. Population Standard Deviation ($\sigma$):
\begin{align*}
    \sigma &= \sqrt{\frac{\sum d_i^2}{n} - \left(\frac{\sum d_i}{n}\right)^2} = \sqrt{\frac{2{,}857}{10} - (4.1)^2} \\
           &= \sqrt{285.7 - 16.81} = \sqrt{268.89} \approx \mathbf{16.40 \text{ mg/dL}}
\end{align*}

3. Sample Standard Deviation ($s$):
\begin{equation*}
    s = \sqrt{\frac{\sum d_i^2 - \frac{(\sum d_i)^2}{n}}{n - 1}} = \sqrt{\frac{2{,}857 - 168.1}{9}} = \sqrt{\frac{2{,}688.9}{9}} \approx \mathbf{17.29 \text{ mg/dL}}
\end{equation*}

\tcbline
\textbf{Part (b): Shift of Origin Method for Decimal Prices:} \\
Let $y_i = X_i - 240$. The transformed observations become:
\begin{center}
    $0.12, \quad 0.13, \quad 0.15, \quad 0.12, \quad 0.17, \quad 0.16, \quad 0.22, \quad 0.21, \quad 0.18, \quad 0.17$
\end{center}
Sum of transformed variables:
\begin{equation*}
    \sum y_i = 1.63, \qquad \sum y_i^2 = 0.2765
\end{equation*}
1. Arithmetic Mean:
\begin{equation*}
    \bar{X} = 240 + \frac{1.63}{10} = \mathbf{240.163}
\end{equation*}

2. Sum of Squared Deviations:
\begin{equation*}
    \sum (X_i - \bar{X})^2 = \sum y_i^2 - \frac{(\sum y_i)^2}{n} = 0.2765 - \frac{(1.63)^2}{10} = 0.2765 - 0.26569 = \mathbf{0.01081}
\end{equation*}

3. Standard Deviation:
\begin{align*}
    \text{Population: } \quad \sigma &= \sqrt{\frac{0.01081}{10}} = \sqrt{0.001081} \approx \mathbf{0.0329} \\
    \text{Sample: } \quad s &= \sqrt{\frac{0.01081}{9}} \approx \mathbf{0.0347}
\end{align*}
\end{example}

\newpage
\lecturedate{29/10/2026}
\begin{example}{Comparative Asset Stability: Reliance vs. JIO Share Prices}{ex_reliance_jio_cv}
\textbf{Problem:} \\
From the consecutive trading price records of Reliance and JIO equity shares furnished below, determine which share is \textbf{more stable in value}:
\begin{center}
\small
\renewcommand{\arraystretch}{1.12}
\begin{tabular}{lcccccccccc}
\toprule
\textbf{Trading Day} & 1 & 2 & 3 & 4 & 5 & 6 & 7 & 8 & 9 & 10 \\
\midrule
\textbf{Reliance (\inr)} & 55 & 54 & 52 & 53 & 56 & 58 & 52 & 50 & 51 & 49 \\
\textbf{JIO (\inr)}      & 108 & 107 & 105 & 105 & 106 & 107 & 104 & 103 & 104 & 101 \\
\bottomrule
\end{tabular}
\end{center}

\tcbline
\textbf{Solution:} \\
To compare relative price stability across two distinct nominal price ranges, we compute the \textbf{Coefficient of Variation ($\CV$)}. A lower $\CV$ signifies smaller relative volatility and superior price stability.

\textbf{Step 1: Statistical Computations for Reliance Shares ($X$):}
\begin{center}
\small
\renewcommand{\arraystretch}{1.12}
\begin{tabularx}{0.75\linewidth}{>{\centering\arraybackslash}p{2.0cm} >{\centering\arraybackslash}p{2.5cm} >{\centering\arraybackslash}p{3.0cm} >{\raggedleft\arraybackslash}X}
\toprule
\textbf{Period ($i$)} & \textbf{Price ($X_i$)} & \textbf{Deviation ($X_i - 53$)} & \textbf{Squared ($(X_i - 53)^2$)} \\
\midrule
1  & 55 & $+2$ & 4 \\
2  & 54 & $+1$ & 1 \\
3  & 52 & $-1$ & 1 \\
4  & 53 & $0$  & 0 \\
5  & 56 & $+3$ & 9 \\
6  & 58 & $+5$ & 25 \\
7  & 52 & $-1$ & 1 \\
8  & 50 & $-3$ & 9 \\
9  & 51 & $-2$ & 4 \\
10 & 49 & $-4$ & 16 \\
\midrule
\textbf{Total ($N=10$)} & $\sum X_i = \mathbf{530}$ & $\sum (X_i - \bar{X}) = 0$ & $\sum (X_i - \bar{X})^2 = \mathbf{70}$ \\
\bottomrule
\end{tabularx}
\end{center}

\begin{align*}
    \bar{X} &= \frac{\sum X_i}{N} = \frac{530}{10} = \mathbf{53.00} \\[0.3em]
    \sigma_{\text{Reliance}} &= \sqrt{\frac{\sum (X_i - \bar{X})^2}{N}} = \sqrt{\frac{70}{10}} = \sqrt{7} \approx \mathbf{2.646} \\[0.3em]
    \CV_{\text{Reliance}} &= \frac{\sigma_{\text{Reliance}}}{\bar{X}} \times 100 = \frac{2.646}{53} \times 100 \approx \mathbf{4.99\%}
\end{align*}

\tcbline
\textbf{Step 2: Statistical Computations for JIO Shares ($Y$):}
\begin{center}
\small
\renewcommand{\arraystretch}{1.12}
\begin{tabularx}{0.75\linewidth}{>{\centering\arraybackslash}p{2.0cm} >{\centering\arraybackslash}p{2.5cm} >{\centering\arraybackslash}p{3.0cm} >{\raggedleft\arraybackslash}X}
\toprule
\textbf{Period ($i$)} & \textbf{Price ($Y_i$)} & \textbf{Deviation ($Y_i - 105$)} & \textbf{Squared ($(Y_i - 105)^2$)} \\
\midrule
1  & 108 & $+3$ & 9 \\
2  & 107 & $+2$ & 4 \\
3  & 105 & $0$  & 0 \\
4  & 105 & $0$  & 0 \\
5  & 106 & $+1$ & 1 \\
6  & 107 & $+2$ & 4 \\
7  & 104 & $-1$ & 1 \\
8  & 103 & $-2$ & 4 \\
9  & 104 & $-1$ & 1 \\
10 & 101 & $-4$ & 16 \\
\midrule
\textbf{Total ($N=10$)} & $\sum Y_i = \mathbf{1{,}050}$ & $\sum (Y_i - \bar{Y}) = 0$ & $\sum (Y_i - \bar{Y})^2 = \mathbf{40}$ \\
\bottomrule
\end{tabularx}
\end{center}

\begin{align*}
    \bar{Y} &= \frac{\sum Y_i}{N} = \frac{1{,}050}{10} = \mathbf{105.00} \\[0.3em]
    \sigma_{\text{JIO}} &= \sqrt{\frac{\sum (Y_i - \bar{Y})^2}{N}} = \sqrt{\frac{40}{10}} = \sqrt{4} = \mathbf{2.000} \\[0.3em]
    \CV_{\text{JIO}} &= \frac{\sigma_{\text{JIO}}}{\bar{Y}} \times 100 = \frac{2.000}{105} \times 100 \approx \mathbf{1.90\%}
\end{align*}

\tcbline
\textbf{Step 3: Comparative Synthesis and Investment Conclusion:}
\begin{center}
\small
\renewcommand{\arraystretch}{1.15}
\begin{tabularx}{\linewidth}{>{\raggedright\arraybackslash}X >{\centering\arraybackslash}p{2.6cm} >{\centering\arraybackslash}p{2.8cm} >{\centering\arraybackslash}p{2.8cm}}
\toprule
\textbf{Security} & \textbf{Mean Price (\inr)} & \textbf{Standard Deviation ($\sigma$)} & \textbf{Coefficient of Variation ($\CV$)} \\
\midrule
Reliance Shares & 53.00  & 2.646 & 4.99\% \\
JIO Shares      & 105.00 & 2.000 & \textbf{1.90\%} \\
\bottomrule
\end{tabularx}
\end{center}

\begin{equation*}
    \CV_{\text{JIO}} \, (1.90\%) < \CV_{\text{Reliance}} \, (4.99\%)
\end{equation*}
\textit{Conclusion:} JIO shares exhibit a substantially lower coefficient of variation ($1.90\%$ vs.~$4.99\%$). Although JIO trades at twice the nominal price of Reliance, its normalized dispersion is over $2.6\times$ smaller. Therefore, \textbf{JIO shares are significantly more stable in value}.
\end{example}

% ==============================================================================
% 4.3 Bivariate Risk, Market Co-movement, and Correlation Mechanics
% ==============================================================================
\newpage
\section{Bivariate Risk, Market Co-movement, and Correlation Mechanics}
\lecturedate{30/10/2026}

\begin{definition}{Covariance and Karl Pearson Correlation Coefficient}{def_covariance_correlation}
In bivariate financial time series, \textbf{covariance} and \textbf{correlation} quantify the directional tendency and linear strength with which two asset return series move in tandem:
\begin{enumerate}[leftmargin=*, topsep=2pt, itemsep=3pt]
    \item \textbf{Covariance ($\Cov(X, Y)$):} The expected value of the cross-product of deviations of two random variables from their respective means:
    \begin{equation}
        \Cov(X, Y) = \frac{1}{n}\sum_{i=1}^{n} (X_i - \bar{X})(Y_i - \bar{Y}) = \frac{\sum X_i Y_i}{n} - \bar{X}\bar{Y}
    \end{equation}
    While covariance establishes the direction of co-movement ($\Cov > 0 \implies$ positive co-movement; $\Cov < 0 \implies$ inverse co-movement), its magnitude is sensitive to measurement units.

    \item \textbf{Karl Pearson Correlation Coefficient ($r_{XY}$):} The standardized, dimensionless normalization of covariance bounded strictly within $[-1, +1]$:
    \begin{equation}
        r_{XY} = \frac{\Cov(X, Y)}{\sigma_X \sigma_Y} = \frac{\sum_{i=1}^{n} (X_i - \bar{X})(Y_i - \bar{Y})}{\sqrt{\sum_{i=1}^{n} (X_i - \bar{X})^2 \sum_{i=1}^{n} (Y_i - \bar{Y})^2}}
    \end{equation}
    \item \textbf{Direct Raw-Score Product-Moment Formula:}
    \begin{equation}
        r_{XY} = \frac{n \sum X_i Y_i - \left(\sum X_i\right)\left(\sum Y_i\right)}{\sqrt{\left[n \sum X_i^2 - \left(\sum X_i\right)^2\right] \left[n \sum Y_i^2 - \left(\sum Y_i\right)^2\right]}}
    \end{equation}
\end{enumerate}
\end{definition}

\lecturedate{30/10/2026}
\begin{example}{Bivariate Co-movement: DLF Stock Returns vs. BSE National Index}{ex_dlf_bse_correlation}
\textbf{Problem:} \\
The monthly percentage return data ($R_m$ and $R_i$) for \textbf{DLF Stock} ($X$) and the \textbf{BSE National Index} ($Y$) across 12 consecutive calendar months are recorded as follows:
\begin{center}
\small
\renewcommand{\arraystretch}{1.15}
\begin{tabular}{lcccccc}
\toprule
\textbf{Month ($i$)} & 1 & 2 & 3 & 4 & 5 & 6 \\
\midrule
\textbf{DLF Return ($X$, \%)} & 9.43 & 0.00 & $-4.31$ & $-18.92$ & $-6.67$ & 26.57 \\
\textbf{BSE Return ($Y$, \%)} & 7.41 & $-5.33$ & $-7.35$ & $-14.64$ & 1.58 & 15.19 \\
\midrule
\textbf{Month ($i$)} & 7 & 8 & 9 & 10 & 11 & 12 \\
\midrule
\textbf{DLF Return ($X$, \%)} & 20.00 & 2.93 & 5.25 & 21.45 & 23.13 & 32.83 \\
\textbf{BSE Return ($Y$, \%)} & 5.11 & 0.76 & $-0.97$ & 10.44 & 17.47 & 20.15 \\
\bottomrule
\end{tabular}
\end{center}
Calculate Karl Pearson's correlation coefficient ($r_{XY}$) and interpret the co-movement profile.

\tcbline
\textbf{Solution:} \\
Let $X$ denote the monthly return of DLF Stock and $Y$ denote the monthly return of the BSE National Index, with sample cardinality $n = 12$.

\textbf{Step 1: Construction of the Product-Moment Calculation Schedule:}
\begin{center}
\small
\renewcommand{\arraystretch}{1.12}
\begin{tabularx}{\linewidth}{>{\centering\arraybackslash}p{1.2cm} >{\raggedleft\arraybackslash}X >{\raggedleft\arraybackslash}X >{\raggedleft\arraybackslash}X >{\raggedleft\arraybackslash}X >{\raggedleft\arraybackslash}X}
\toprule
\textbf{Month} & \textbf{DLF ($X$)} & \textbf{BSE ($Y$)} & \textbf{$X^2$} & \textbf{$Y^2$} & \textbf{$XY$} \\
\midrule
1  & 9.43   & 7.41   & 88.9249   & 54.9081   & 69.8763 \\
2  & 0.00   & $-5.33$ & 0.0000    & 28.4089   & 0.0000 \\
3  & $-4.31$ & $-7.35$ & 18.5761   & 54.0225   & 31.6785 \\
4  & $-18.92$ & $-14.64$ & 357.9664 & 214.3296  & 276.9888 \\
5  & $-6.67$ & 1.58   & 44.4889   & 2.4964    & $-10.5386$ \\
6  & 26.57  & 15.19  & 705.9649  & 230.7361  & 403.5983 \\
7  & 20.00  & 5.11   & 400.0000  & 26.1121   & 102.2000 \\
8  & 2.93   & 0.76   & 8.5849    & 0.5776    & 2.2268 \\
9  & 5.25   & $-0.97$ & 27.5625   & 0.9409    & $-5.0925$ \\
10 & 21.45  & 10.44  & 460.1025  & 108.9936  & 223.9380 \\
11 & 23.13  & 17.47  & 534.9969  & 305.2009  & 404.0811 \\
12 & 32.83  & 20.15  & 1077.8089 & 406.0225  & 661.5245 \\
\midrule
\textbf{Total} & \textbf{111.69} & \textbf{49.82} & \textbf{3,724.9769} & \textbf{1,432.7492} & \textbf{2,160.4812} \\
\bottomrule
\end{tabularx}
\end{center}

\tcbline
\textbf{Step 2: Substitution into the Direct Pearson Formulation:}
\begin{align*}
    n \sum XY - \left(\sum X\right)\left(\sum Y\right) &= 12(2{,}160.4812) - (111.69)(49.82) \\
                                                      &= 25{,}925.7744 - 5{,}564.3958 = \mathbf{20{,}361.3786} \\[0.5em]
    n \sum X^2 - \left(\sum X\right)^2 &= 12(3{,}724.9769) - (111.69)^2 \\
                                      &= 44{,}699.7228 - 12{,}474.6561 = \mathbf{32{,}225.0667} \\[0.5em]
    n \sum Y^2 - \left(\sum Y\right)^2 &= 12(1{,}432.7492) - (49.82)^2 \\
                                      &= 17{,}192.9904 - 2{,}482.0324 = \mathbf{14{,}710.9580}
\end{align*}

Evaluating the denominator product and its square root:
\begin{equation*}
    \sqrt{32{,}225.0667 \times 14{,}710.9580} = \sqrt{474{,}061{,}699.7} \approx \mathbf{21{,}772.9558}
\end{equation*}

Computing the correlation coefficient $r_{XY}$:
\begin{equation*}
    r_{XY} = \frac{20{,}361.3786}{21{,}772.9558} \approx \mathbf{+0.9352} \quad (\approx \mathbf{+0.935})
\end{equation*}

\tcbline
\textbf{Step 3: Financial Econometric Interpretation:}
\begin{itemize}[leftmargin=*, noitemsep]
    \item \textbf{Strength of Association:} The estimated correlation $r \approx +0.935$ reveals an exceptionally strong, statistically robust positive linear association between DLF stock returns and the benchmark BSE National Index.
    \item \textbf{Systematic Market Risk Exposure:} DLF shares are driven overwhelmingly by macroeconomic market factors ($\approx 87.5\%$ of its variance, $R^2 = r^2 = 0.8746$, is explained by broad market index co-movements).
    \item \textbf{Portfolio Beta Implication:} Because $r_{XY} > 0.90$, holding DLF stock provides minimal diversification benefit against broad market pullbacks, demonstrating pronounced systematic market beta risk.
\end{itemize}
\end{example}

% ==============================================================================
% 4.4 Sharpe's Single-Index Optimal Portfolio Construction
% ==============================================================================
\newpage
\section{Sharpe's Single-Index Optimal Portfolio Construction}
\lecturedate{30/10/2026}

\subsection{Theoretical Foundations of the Single-Index Model}
In Harry Markowitz's classical mean-variance portfolio theory, optimizing an $N$-asset portfolio requires estimating $N$ expected returns, $N$ asset variances, and $N(N-1)/2$ pairwise covariances—creating severe computational intractability and parameter error for large investment universes ($N = 100 \implies 4{,}950$ covariances). 

To overcome this curse of dimensionality, William Sharpe (1963) developed the \textbf{Single-Index Model (SIM)}. The SIM posits that the co-movement between individual securities arises entirely from a common macroeconomic factor—represented by the broad market index return ($R_m$):
\begin{equation}
    R_i = \alpha_i + \beta_i R_m + e_i
\end{equation}
subject to the standard Gauss--Markov linear index assumptions:
\begin{align}
    E[e_i] &= 0 \quad \text{for all } i, \\[0.3em]
    \Var(e_i) &= \sigma_{ei}^2 \quad (\text{constant unsystematic / residual variance}), \\[0.3em]
    \Cov(e_i, e_j) &= 0 \quad \text{for all } i \neq j \quad (\text{zero residual cross-correlation}), \\[0.3em]
    \Cov(e_i, R_m) &= 0 \quad \text{for all } i \quad (\text{residuals orthogonal to market return}).
\end{align}

Under these assumptions, the total risk of any security $i$ decomposes into two mutually independent, orthogonal components:
\begin{equation}
    \sigma_i^2 = \underbrace{\beta_i^2 \sigma_m^2}_{\text{Systematic Risk}} + \underbrace{\sigma_{ei}^2}_{\text{Unsystematic (Firm-Specific) Risk}}
\end{equation}
and the covariance between any two distinct assets $i$ and $j$ ($i \neq j$) depends solely on their joint systematic exposure to the market index:
\begin{equation}
    \Cov(R_i, R_j) = \beta_i \beta_j \sigma_m^2
\end{equation}

\subsection{Sharpe Ratio Maximization and Optimal Cut-off Rate \texorpdfstring{$C^*$}{C*}}
The investor's objective is to construct an optimal tangent portfolio of risky assets that maximizes the \textbf{Sharpe ratio} ($\theta_p$), defined as the slope of the Capital Allocation Line (CAL):
\begin{equation}
    \max_{\{w_i\}} \theta_p = \frac{E[R_p] - R_f}{\sigma_p} \quad \text{subject to} \quad \sum_{i=1}^N w_i = 1, \quad w_i \ge 0 \quad (\text{no short sales})
\end{equation}

Elton, Gruber, and Padberg (1976) proved that maximizing $\theta_p$ under the Single-Index Model yields a unique optimal cut-off acceptance threshold, denoted by $C^*$. A security qualifies for inclusion in the optimal portfolio if and only if its excess return-to-beta ratio exceeds this critical cut-off rate:
\begin{equation}
    Z_i > C^* \iff \text{Include Security } i \text{ in the Optimal Portfolio}
\end{equation}
Securities with $Z_i \le C^*$ are rejected because their inclusion would dilute the portfolio's Sharpe ratio.

\subsection{Algorithmic Optimization Pipeline}
The mathematical process of constructing Sharpe's optimal portfolio follows a four-stage systematic pipeline, illustrated in Figure~\ref{fig:sharpe_pipeline}.

\begin{figure}[H]
\centering
\begin{tikzpicture}[>=stealth, font=\small]
    \tikzset{
        stagecard/.style={
            draw=navyblue!80!black,
            fill=navyblue!4,
            rounded corners=4pt,
            line width=1.1pt,
            text width=3.2cm,
            minimum height=3.8cm,
            align=center,
            inner sep=4pt
        },
        transarrow/.style={
            ->,
            line width=1.3pt,
            navyblue,
            shorten >=2pt,
            shorten <=2pt
        }
    }

    % Stage 1 Card
    \node[stagecard] (card1) at (0, 0) {
        {\bfseries\color{navyblue} Stage 1: Security Ranking}\\[0.35em]
        \scriptsize Compute excess return to beta ratio:\\[0.25em]
        $Z_i = \dfrac{R_i - R_f}{\beta_i}$\\[0.4em]
        Rank all candidate assets in strictly descending order:\\[0.2em]
        $Z_1 \ge Z_2 \ge \dots \ge Z_N$
    };

    % Stage 2 Card
    \node[stagecard] (card2) at (3.9, 0) {
        {\bfseries\color{teal!80!black} Stage 2: Variance Factors}\\[0.35em]
        \scriptsize Compute variance-weighted factors:\\[0.2em]
        $A_i = \dfrac{(R_i - R_f)\beta_i}{\sigma_{ei}^2}$\\[0.3em]
        $B_i = \dfrac{\beta_i^2}{\sigma_{ei}^2}$\\[0.3em]
        Accumulate sums $\sum A_i, \sum B_i$
    };

    % Stage 3 Card
    \node[stagecard] (card3) at (7.8, 0) {
        {\bfseries\color{forestgreen!80!black} Stage 3: Cut-off Rate $C^*$}\\[0.35em]
        \scriptsize Evaluate cumulative cut-off rate:\\[0.2em]
        $C_i = \dfrac{\sigma_m^2 \sum A_i}{1 + \sigma_m^2 \sum B_i}$\\[0.35em]
        Identify acceptance boundary:\\[0.2em]
        $Z_i > C_i \implies C^* = \max C_i$
    };

    % Stage 4 Card
    \node[stagecard] (card4) at (11.7, 0) {
        {\bfseries\color{crimson!80!black} Stage 4: Capital Weights}\\[0.35em]
        \scriptsize Compute relative unscaled weights:\\[0.2em]
        $Y_i = \dfrac{\beta_i}{\sigma_{ei}^2}(Z_i - C^*)$\\[0.35em]
        Normalize to percentage capital allocations:\\[0.2em]
        $w_i = \dfrac{Y_i}{\sum Y_i}$
    };

    % Transition Connectors with ample clearance
    \draw[transarrow] (card1) -- (card2);
    \draw[transarrow] (card2) -- (card3);
    \draw[transarrow] (card3) -- (card4);
\end{tikzpicture}
\caption{Systematic Four-Stage Architecture of Sharpe's Single-Index Optimal Portfolio Selection.}
\label{fig:sharpe_pipeline}
\end{figure}

\begin{definition}{Sharpe's Optimal Portfolio Selection Formulations}{def_sharpe_sim_opt}
Let $R_f$ denote the riskless rate of return and $\sigma_m^2$ represent the variance of the macroeconomic market index. For an investment universe of $N$ candidate risky securities:
\begin{enumerate}[leftmargin=*, topsep=2pt, itemsep=3pt]
    \item \textbf{Excess Return-to-Beta Ratio (Ranking Metric):}
    \begin{equation}
        Z_i = \frac{R_i - R_f}{\beta_i}
    \end{equation}
    Securities are ranked from highest to lowest: $Z_1 \ge Z_2 \ge \dots \ge Z_N$.

    \item \textbf{Individual Variance-Weighted Factors:}
    \begin{equation}
        A_i = \frac{(R_i - R_f)\beta_i}{\sigma_{ei}^2}, \qquad B_i = \frac{\beta_i^2}{\sigma_{ei}^2}
    \end{equation}

    \item \textbf{Cumulative Cut-off Rate ($C_i$):}
    For a portfolio containing the top $i$ ranked securities:
    \begin{equation}
        C_i = \frac{\sigma_m^2 \sum_{j=1}^i A_j}{1 + \sigma_m^2 \sum_{j=1}^i B_j} = \frac{\sigma_m^2 \sum_{j=1}^i \frac{(R_j - R_f)\beta_j}{\sigma_{ej}^2}}{1 + \sigma_m^2 \sum_{j=1}^i \frac{\beta_j^2}{\sigma_{ej}^2}}
    \end{equation}

    \item \textbf{Asset Selection Condition and Optimal Cut-off Point ($C^*$):}
    A security $i$ is included in the optimal portfolio if and only if:
    \begin{equation}
        Z_i > C_i
    \end{equation}
    The optimal cut-off threshold is $C^* = C_k$, where $k$ is the last asset index satisfying $Z_k > C_k$. At asset $k+1$, $Z_{k+1} \le C_k$, terminating the search.

    \item \textbf{Unscaled Weights and Normalized Portfolio Allocations:}
    For each selected security $i \in \{1, 2, \dots, k\}$:
    \begin{equation}
        Y_i = \frac{\beta_i}{\sigma_{ei}^2} (Z_i - C^*), \qquad w_i = \frac{Y_i}{\sum_{j=1}^k Y_j} \quad \text{such that} \quad \sum_{i=1}^k w_i = 1
    \end{equation}
\end{enumerate}
\end{definition}

\begin{tcolorbox}[colback=lightgray!15, colframe=navyblue!60, title=\textbf{Pedagogical Note on Notation Conventions}]
\small
In the classroom lecture notes, the symbol $Z_i$ denotes the excess return-to-beta ranking ratio ($\frac{R_i - R_f}{\beta_i}$) and $Y_i$ denotes the unscaled relative asset weights. In standard textbook literature (e.g., Elton, Gruber, Brown, and Goetzmann, \textit{Modern Portfolio Theory and Investment Analysis}), the excess return-to-beta ratio is written directly as $\frac{R_i - R_f}{\beta_i}$, while the unscaled weights are denoted by $z_i$. Both mathematical systems are identical in logic, intermediate values, and final capital allocations ($w_i$).
\end{tcolorbox}

\begin{example}{Optimal Portfolio Construction via Single-Index Model (Jagan's Portfolio)}{ex_jagan_sharpe_opt}
\lecturedate{30/10/2026}
\textbf{Problem:}\\[0.2em]
Jagan is constructing an optimal investment portfolio from a candidate universe of five corporate stocks. The market variance is $\sigma_m^2 = 10\%$, and the prevailing riskless rate of return is $R_f = 5\%$. The expected returns, beta coefficients, and unsystematic residual risk levels of the five companies are presented in the initial portfolio data table below.
\begin{enumerate}[leftmargin=*, noitemsep]
    \item Rank the candidate stocks by their excess return-to-beta ratio ($Z_i$).
    \item Construct the cumulative cut-off schedule ($C_i$) and determine the optimal cut-off point ($C^*$).
    \item Identify the specific companies selected for inclusion in the optimum portfolio.
    \item Calculate the relative unscaled weights ($Y_i$) and final percentage capital allocations ($w_i$).
    \item Perform a post-optimization portfolio synthesis check: compute the portfolio's expected return $E[R_p]$, systematic beta $\beta_p$, unsystematic risk $\sigma_{ep}^2$, total variance $\sigma_p^2$, and Sharpe ratio $\theta_p$.
\end{enumerate}

\begin{center}
\small
\begin{minipage}{0.88\linewidth}
\centering
\textbf{Table 1: Given Initial Portfolio Inputs}\\[0.4em]
\renewcommand{\arraystretch}{1.15}
\begin{tabularx}{\linewidth}{>{\centering\arraybackslash}p{2.5cm} >{\centering\arraybackslash}X >{\centering\arraybackslash}X >{\centering\arraybackslash}X}
\toprule
\textbf{Company ($i$)} & \textbf{Return ($R_i, \%$)} & \textbf{Beta ($\beta_i$)} & \textbf{Unsystematic Risk ($\sigma_{ei}^2$)} \\
\midrule
1 & 20 & 1.0 & 20 \\
2 & 22 & 1.5 & 10 \\
3 & 10 & 0.5 & 30 \\
4 & 24 & 2.0 & 20 \\
5 & 12 & 1.0 & 30 \\
\bottomrule
\end{tabularx}
\end{minipage}
\end{center}

\tcbline
\textbf{Solution:}

\textbf{Step 1: Calculation of Excess Return to Beta ($Z_i$) and Security Ranking:}
Using the formula $Z_i = \frac{R_i - R_f}{\beta_i}$ with $R_f = 5\%$:
\begin{itemize}[leftmargin=*, noitemsep]
    \item Company 1: $Z_1 = \frac{20 - 5}{1.0} = \frac{15}{1.0} = 15.00$
    \item Company 2: $Z_2 = \frac{22 - 5}{1.5} = \frac{17}{1.5} = 11.33$
    \item Company 3: $Z_3 = \frac{10 - 5}{0.5} = \frac{5}{0.5} = 10.00$
    \item Company 4: $Z_4 = \frac{24 - 5}{2.0} = \frac{19}{2.0} = 9.50$
    \item Company 5: $Z_5 = \frac{12 - 5}{1.0} = \frac{7}{1.0} = 7.00$
\end{itemize}

\begin{center}
\small
\begin{minipage}{0.92\linewidth}
\centering
\textbf{Table 2: Step 1 --- Excess Return to Beta ($Z_i$) Ranking Schedule}\\[0.4em]
\renewcommand{\arraystretch}{1.15}
\begin{tabularx}{\linewidth}{>{\centering\arraybackslash}p{2.0cm} >{\centering\arraybackslash}X >{\centering\arraybackslash}X >{\centering\arraybackslash}X >{\centering\arraybackslash}X >{\centering\arraybackslash}p{1.8cm}}
\toprule
\textbf{Company} & \textbf{Return ($R_i$)} & \textbf{$R_i - R_f$} & \textbf{Beta ($\beta_i$)} & \textbf{$Z_i$ Ratio} & \textbf{Rank} \\
\midrule
1 & 20 & 15 & 1.0 & \textbf{15.00} & \textbf{1} \\
2 & 22 & 17 & 1.5 & \textbf{11.33} & \textbf{2} \\
3 & 10 & 5  & 0.5 & \textbf{10.00} & \textbf{3} \\
4 & 24 & 19 & 2.0 & \textbf{9.50}  & \textbf{4} \\
5 & 12 & 7  & 1.0 & \textbf{7.00}  & \textbf{5} \\
\bottomrule
\end{tabularx}
\end{minipage}
\end{center}

\tcbline
\textbf{Step 2: Calculation of the Cumulative Cut-off Schedule ($C_i$):}
For each ranked company, evaluate:
\begin{equation*}
    A_i = \frac{(R_i - R_f)\beta_i}{\sigma_{ei}^2}, \qquad B_i = \frac{\beta_i^2}{\sigma_{ei}^2}, \qquad C_i = \frac{\sigma_m^2 \sum_{j=1}^i A_j}{1 + \sigma_m^2 \sum_{j=1}^i B_j} \quad (\text{with } \sigma_m^2 = 10)
\end{equation*}

\begin{center}
\footnotesize
\begin{minipage}{\linewidth}
\centering
\textbf{Table 3: Step 2 --- Cumulative Cut-off Rate ($C_i$) Computation Schedule}\\[0.4em]
\setlength{\tabcolsep}{2.5pt}
\renewcommand{\arraystretch}{1.2}
\begin{tabularx}{\linewidth}{>{\centering\arraybackslash}p{0.9cm} >{\centering\arraybackslash}p{0.6cm} >{\centering\arraybackslash}p{0.8cm} >{\centering\arraybackslash}p{1.7cm} >{\centering\arraybackslash}p{0.8cm} >{\centering\arraybackslash}p{1.4cm} >{\centering\arraybackslash}p{0.8cm} >{\centering\arraybackslash}X >{\centering\arraybackslash}p{1.0cm}}
\toprule
\textbf{Rank} & \textbf{Co.} & \textbf{$Z_i$} & \textbf{$A_i = \frac{(R_i - R_f)\beta_i}{\sigma_{ei}^2}$} & \textbf{$\sum A_i$} & \textbf{$B_i = \frac{\beta_i^2}{\sigma_{ei}^2}$} & \textbf{$\sum B_i$} & \textbf{$C_i$ Calculation} & \textbf{$C_i$} \\
\midrule
1 & 1 & 15.00 & $\frac{15 \times 1.0}{20} = 0.750$ & 0.750 & $\frac{1.0^2}{20} = 0.050$ & 0.050 & $\frac{10 \times 0.750}{1 + 10 \times 0.050} = \frac{7.500}{1.500}$ & \textbf{5.000} \\[0.4em]
2 & 2 & 11.33 & $\frac{17 \times 1.5}{10} = 2.550$ & 3.300 & $\frac{1.5^2}{10} = 0.225$ & 0.275 & $\frac{10 \times 3.300}{1 + 10 \times 0.275} = \frac{33.000}{3.750}$ & \textbf{8.800} \\[0.4em]
3 & 3 & 10.00 & $\frac{5 \times 0.5}{30} = 0.083$  & 3.383 & $\frac{0.5^2}{30} = 0.008$ & 0.283 & $\frac{10 \times 3.383}{1 + 10 \times 0.283} = \frac{33.833}{3.833}$ & \textbf{8.826} \\[0.4em]
4 & 4 & 9.50  & $\frac{19 \times 2.0}{20} = 1.900$ & 5.283 & $\frac{2.0^2}{20} = 0.200$ & 0.483 & $\frac{10 \times 5.283}{1 + 10 \times 0.483} = \frac{52.833}{5.833}$ & \textbf{9.057} \\[0.4em]
5 & 5 & 7.00  & $\frac{7 \times 1.0}{30} = 0.233$  & 5.516 & $\frac{1.0^2}{30} = 0.033$ & 0.516 & $\frac{10 \times 5.516}{1 + 10 \times 0.516} = \frac{55.167}{6.167}$ & \textit{8.946} \\
\bottomrule
\end{tabularx}
\end{minipage}
\end{center}

\textbf{Determining the Optimal Cut-off Point ($C^*$):}
A security is included in the portfolio as long as its excess return-to-beta ratio exceeds the cumulative cut-off rate ($Z_i > C_i$):
\begin{itemize}[leftmargin=*, noitemsep]
    \item Company 1: $Z_1 (15.00) > C_1 (5.000) \implies$ Include
    \item Company 2: $Z_2 (11.33) > C_2 (8.800) \implies$ Include
    \item Company 3: $Z_3 (10.00) > C_3 (8.826) \implies$ Include
    \item Company 4: $Z_4 (9.50) > C_4 (9.057) \implies$ Include
    \item Company 5: $Z_5 (7.00) < C_4 (9.057) \implies$ \textbf{Exclude / Stop}
\end{itemize}
Because $Z_5 = 7.00$ falls below the prevailing cut-off rate $C_4 = 9.057$, Company 5 is excluded, and the optimization process terminates.

\textbf{Conclusion on Security Selection:}
\begin{equation*}
    \text{Optimal Cut-off Point: } \mathbf{C^* = 9.057}
\end{equation*}
The optimum portfolio consists strictly of \textbf{Companies 1, 2, 3, and 4}.

\tcbline
\textbf{Step 3: Calculation of Optimum Portfolio Weights ($w_i$):}
For each selected asset $i \in \{1, 2, 3, 4\}$, the unscaled relative weight $Y_i$ is computed as:
\begin{equation*}
    Y_i = \frac{\beta_i}{\sigma_{ei}^2} (Z_i - C^*) = \frac{\beta_i}{\sigma_{ei}^2} (Z_i - 9.057)
\end{equation*}
and the final percentage capital allocation is obtained via normalization: $w_i = \frac{Y_i}{\sum_{j=1}^4 Y_j}$.

\begin{center}
\small
\begin{minipage}{\linewidth}
\centering
\textbf{Table 4: Step 3 --- Optimum Asset Weights and Percentage Capital Allocation}\\[0.4em]
\setlength{\tabcolsep}{3pt}
\renewcommand{\arraystretch}{1.2}
\begin{tabularx}{\linewidth}{>{\centering\arraybackslash}p{1.7cm} >{\centering\arraybackslash}p{0.9cm} >{\centering\arraybackslash}p{1.0cm} >{\centering\arraybackslash}p{0.9cm} >{\centering\arraybackslash}p{1.6cm} >{\raggedleft\arraybackslash}X >{\centering\arraybackslash}p{2.2cm}}
\toprule
\textbf{Company} & \textbf{Beta ($\beta_i$)} & \textbf{$\sigma_{ei}^2$} & \textbf{$Z_i$} & \textbf{$Z_i - C^*$} & \textbf{Relative Weight ($Y_i$)} & \textbf{Final Weight ($w_i$)} \\
\midrule
1 & 1.0 & 20 & 15.00 & $15.00 - 9.057 = 5.943$ & $\frac{1.0}{20} \times 5.943 = 0.29715$ & $\frac{0.29715}{0.69857} \approx \mathbf{42.54\%}$ \\[0.3em]
2 & 1.5 & 10 & 11.33 & $11.33 - 9.057 = 2.276$ & $\frac{1.5}{10} \times 2.276 = 0.34140$ & $\frac{0.34140}{0.69857} \approx \mathbf{48.87\%}$ \\[0.3em]
3 & 0.5 & 30 & 10.00 & $10.00 - 9.057 = 0.943$ & $\frac{0.5}{30} \times 0.943 = 0.01572$ & $\frac{0.01572}{0.69857} \approx \mathbf{2.25\%}$ \\[0.3em]
4 & 2.0 & 20 & 9.50  & $9.50 - 9.057 = 0.443$  & $\frac{2.0}{20} \times 0.443 = 0.04430$ & $\frac{0.04430}{0.69857} \approx \mathbf{6.34\%}$ \\
\midrule
\textbf{Total} & — & — & — & — & $\sum Y_i = \mathbf{0.69857}$ & $\mathbf{100.00\%}$ \\
\bottomrule
\end{tabularx}
\end{minipage}
\end{center}

\tcbline
\textbf{Step 4: Post-Optimization Portfolio Synthesis Check:}
\begin{enumerate}[leftmargin=*, itemsep=3pt]
    \item \textbf{Expected Portfolio Return ($E[R_p]$):}
    \begin{align*}
        E[R_p] &= \sum_{i=1}^4 w_i R_i = 0.4254(20) + 0.4887(22) + 0.0225(10) + 0.0634(24) \\
               &= 8.508\% + 10.751\% + 0.225\% + 1.522\% = \mathbf{21.01\%}
    \end{align*}

    \item \textbf{Portfolio Beta Sensitivity ($\beta_p$):}
    \begin{align*}
        \beta_p &= \sum_{i=1}^4 w_i \beta_i = 0.4254(1.0) + 0.4887(1.5) + 0.0225(0.5) + 0.0634(2.0) \\
                &= 0.4254 + 0.7331 + 0.0113 + 0.1268 = \mathbf{1.2966} \approx \mathbf{1.30}
    \end{align*}

    \item \textbf{Portfolio Unsystematic Variance ($\sigma_{ep}^2$):}
    \begin{align*}
        \sigma_{ep}^2 &= \sum_{i=1}^4 w_i^2 \sigma_{ei}^2 = (0.4254)^2(20) + (0.4887)^2(10) + (0.0225)^2(30) + (0.0634)^2(20) \\
                      &= 3.6193 + 2.3883 + 0.0152 + 0.0804 = \mathbf{6.1032}
    \end{align*}

    \item \textbf{Total Optimal Portfolio Variance and Standard Deviation ($\sigma_p^2, \sigma_p$):}
    \begin{align*}
        \sigma_p^2 &= \beta_p^2 \sigma_m^2 + \sigma_{ep}^2 = (1.2966)^2(10) + 6.1032 = 16.8117 + 6.1032 = \mathbf{22.9149} \\[0.3em]
        \sigma_p   &= \sqrt{22.9149} \approx \mathbf{4.787\%}
    \end{align*}

    \item \textbf{Optimal Tangency Portfolio Sharpe Ratio ($\theta_p$):}
    \begin{equation*}
        \theta_p = \frac{E[R_p] - R_f}{\sigma_p} = \frac{21.006\% - 5.000\%}{4.787\%} = \frac{16.006}{4.787} \approx \mathbf{3.3436}
    \end{equation*}
    The assembled portfolio achieves a high Sharpe ratio of $3.34$, demonstrating that Sharpe's single-index cut-off methodology successfully balances high expected return against systematic and residual risk.
\end{enumerate}
\end{example}

% ==============================================================================
% 4.5 Foreign Exchange Exposure and Currency Risk Foundations
% ==============================================================================
\newpage
\section{Foreign Exchange Exposure and Currency Risk Foundations}
\lecturedate{30/10/2026}

\begin{definition}{Currency Exposure versus Exchange Rate Risk}{def_forex_exposure}
In multinational financial management and quantitative corporate finance, international cash flows and firm valuation are subject to foreign exchange rate volatility:
\begin{enumerate}[leftmargin=*, topsep=2pt, itemsep=3pt]
    \item \textbf{Foreign Exchange Exposure (Sensitivity):} The degree of structural sensitivity or responsiveness with which a corporation's operating cash flows, domestic asset values, or stock returns are affected by unpredictable exchange rate fluctuations.
    Econometrically, currency exposure is estimated via linear regression modeling:
    \begin{equation}
        R_{i, t} = \alpha_i + \beta_{\text{FX}} \left(\frac{\Delta S_t}{S_{t-1}}\right) + \epsilon_t
    \end{equation}
    where $\frac{\Delta S_t}{S_{t-1}}$ represents the percentage change in the spot exchange rate, and the regression slope $\beta_{\text{FX}} = \frac{\Cov(R_i, \Delta S/S)}{\Var(\Delta S/S)}$ represents the firm's \textbf{foreign exchange exposure coefficient}.

    \item \textbf{Exchange Rate Risk (Variance / Volatility):} The actual variance, dispersion, or standard deviation in the firm's equity value or cash flows attributable directly to exchange rate movements:
    \begin{equation}
        \Var(R_{i}) = \beta_{\text{FX}}^2 \sigma_{\text{FX}}^2 + \sigma_\epsilon^2
    \end{equation}
    While exposure reflects the structural slope ($\beta_{\text{FX}}$), risk combines exposure with the underlying market volatility ($\sigma_{\text{FX}}^2$).

    \item \textbf{FOREX Exposure (Cross-Border Capital Activities):} The overall enterprise vulnerability arising from commercial trade, offshore investments, cross-border liabilities, or debt financing denominated in currencies other than the home domestic reporting currency.
\end{enumerate}
\end{definition}

\subsection{Tripartite Taxonomy of Foreign Exchange Exposure}
Multinational enterprises classify currency exposures into three operational categories:

\begin{center}
\small
\renewcommand{\arraystretch}{1.2}
\begin{tabularx}{\linewidth}{>{\raggedright\arraybackslash}p{3.2cm} >{\raggedright\arraybackslash}p{3.2cm} >{\raggedright\arraybackslash}X >{\raggedright\arraybackslash}p{3.0cm}}
\toprule
\textbf{Exposure Type} & \textbf{Temporal Horizon} & \textbf{Economic Impact / Focus} & \textbf{Primary Hedging Tools} \\
\midrule
\textbf{Transaction Exposure} & Short-term ($\le 1$ year) & Contracted foreign currency receivables, payables, and debt obligations with fixed settlement dates. & Forward contracts, futures, money market hedges, currency options. \\
\addlinespace
\textbf{Translation Exposure} & Periodic (Quarterly / Annual reporting) & Paper accounting volatility arising when consolidating foreign subsidiary balance sheets into parent currency. & Balance sheet hedging, matching local currency assets with liabilities. \\
\addlinespace
\textbf{Economic / Operating Exposure} & Long-term strategic horizon (Multi-year) & Uncontracted future operating cash flows and competitive market pricing power altered by structural exchange rate realignments. & Operational diversification, production relocation, multi-currency financing. \\
\bottomrule
\end{tabularx}
\end{center}
